\begin{enumerate}
\item Đường tròn ngoại tiếp tam giác là đường tròn đi qua $3$ đỉnh (Dạng 2). Tâm là giao $3$ đường trung trực.
\item Đường tròn nội tiếp tam giác tiếp xúc $3$ cạnh.
\item Kết hợp các yếu tố hình học phẳng (góc, khoảng cách) đưa vào hệ tọa độ.
\end{enumerate}
Lưu ý: Tọa độ tâm đường tròn ngoại tiếp có thể tìm nhanh bằng cách gọi $I(x; y)$ rồi giải hệ $IA = IB$ và $IA = IC$.
